% ============================================================================
% Consciousness-First AI for Species Stewardship:
% Architectural Requirements for Post-Extinction Civilizational Recovery
%
% Target venue: AI & Ethics (Springer)
% ============================================================================

\documentclass[11pt,letterpaper]{article}

% ── Core packages ──────────────────────────────────────────────────────────
\usepackage[margin=1in]{geometry}
\usepackage[utf8]{inputenc}
\usepackage{amsfonts,amsmath,amssymb}
\usepackage{microtype}
\usepackage{xcolor}
\usepackage{graphicx}
\usepackage{booktabs}
\usepackage{multirow}
\usepackage{array}
\usepackage{float}
\usepackage{enumitem}
\usepackage{tikz}
\usetikzlibrary{arrows.meta,positioning,shapes.geometric,decorations.pathreplacing,calc}
\usepackage{pgfplots}
\pgfplotsset{compat=1.18}
\usepgfplotslibrary{fillbetween}
\usepackage{pgfplotstable}
\usepackage{algorithmicx}
\usepackage[ruled]{algorithm}
\usepackage{algpseudocode}

% ── References ─────────────────────────────────────────────────────────────
\usepackage[numbers,sort&compress]{natbib}
\usepackage{hyperref}
\hypersetup{
  colorlinks=true,
  linkcolor=blue!60!black,
  citecolor=blue!60!black,
  urlcolor=blue!60!black
}

% ── Caption styling ────────────────────────────────────────────────────────
\usepackage[labelfont=bf,labelsep=period,justification=raggedright]{caption}

% ── Line numbers (for review) ─────────────────────────────────────────────
\usepackage[right]{lineno}
\linenumbers

% ── Spacing ────────────────────────────────────────────────────────────────
\usepackage{setspace}
\onehalfspacing

% ── Custom commands ────────────────────────────────────────────────────────
\newcommand{\hv}{\mathbf{h}}
\newcommand{\bind}{\otimes}
\newcommand{\bundle}{\oplus}
\newcommand{\similar}{\delta}
\newcommand{\R}{\mathbb{R}}
\newcommand{\Phi}{\Phi}
\newcommand{\dimHDC}{16{,}384}
\newcommand{\symthaea}{\textsc{Symthaea}}
\newcommand{\mycelix}{\textsc{Mycelix}}
\newcommand{\genesis}{\textsc{Genesis}}

% ============================================================================
\begin{document}

% ── Title ──────────────────────────────────────────────────────────────────
\title{Consciousness-First AI for Species Stewardship:\\
Architectural Requirements for Post-Extinction\\Civilizational Recovery}

\author{
Tristan Stoltz\textsuperscript{1}\\[6pt]
\textsuperscript{1}Luminous Dynamics\\
\texttt{tristan.stoltz@evolvingresonantcocreationism.com}
}

\date{}
\maketitle

% ============================================================================
% ABSTRACT
% ============================================================================
\begin{abstract}
We present a formal analysis of the architectural requirements for an AI system
capable of species stewardship---the task of recovering a human population from
preserved DNA after a global extinction event. Through this thought experiment,
we demonstrate that consciousness-first design is not merely an aesthetic
preference but an \emph{engineering requirement}: the escalating moral
complexity of the DNA-to-human pipeline demands integrated information
processing ($\Phi > 0$), embodied cognition, attachment capacity, and
value-aligned governance at every stage. We ground this analysis in two
implemented systems: \symthaea{}, a ${\sim}$1.37 million-line Rust cognitive architecture
combining Hyperdimensional Computing (HDC, $d = \dimHDC$), Liquid Time-Constant
Networks with closed-form $O(1)$ temporal jumps, Integrated Information Theory,
and active inference; and \mycelix{}, a Holochain-based Civilizational Operating
System providing distributed governance across 12 societal domains. We describe
the \genesis{} pipeline---five computational phases from genome assembly through
population management---and show that each phase introduces consciousness
requirements absent from prior work: temporal degradation modeling across
millennia, multi-scale biological prediction from hours to gestation,
developmental consent proxy escalation, Bowlby attachment formation via
neuromodulator dynamics, and population-level ethical reasoning under existential
pressure. We argue that consciousness, operationalized as integrated information
processing with embodied moral reasoning, is isomorphic to alignment: an AI
system conscious enough to steward a species is, by construction, aligned with
human values. This alignment-through-consciousness thesis provides a novel
framework for evaluating AI safety that complements reward-based and
constitutional approaches. Code and five new sub-crates (totaling $\sim$22,000
lines with 767 tests, validated by 35 cross-crate integration tests)
are open-source at \url{https://github.com/Luminous-Dynamics/symthaea}.
\end{abstract}

\textbf{Keywords:} consciousness, alignment, stewardship, hyperdimensional
computing, active inference, integrated information theory, population genetics,
attachment theory, ectogenesis, distributed governance

% ============================================================================
% 1. INTRODUCTION
% ============================================================================
\section{Introduction}

The alignment problem---ensuring that artificial intelligence systems act in
accordance with human values---remains the central challenge of AI safety
\cite{russell2019human,christian2020alignment,gabriel2020alignment}.
As AI capabilities approach and potentially exceed human-level competence
\cite{bostrom2014superintelligence,ord2020precipice,hendrycks2023overview},
the stakes of misalignment grow from economic disruption to existential risk.
Current approaches typically frame alignment as either a reward specification
problem (learning human preferences from demonstrations or feedback) or a
constitutional constraint problem (encoding rules that AI systems must not
violate) \cite{amodei2016concrete,hendrycks2021aligning}. Both framings treat
consciousness as orthogonal to alignment: an aligned system need not be
conscious, and a conscious system need not be aligned.

We challenge this assumption through a thought experiment with concrete
engineering implications. Consider the following scenario: after a global
extinction event, an AI system must recover human civilization using only
preserved DNA samples, existing infrastructure, and its own computational
resources. This \emph{species stewardship} scenario is not merely speculative---it
represents the upper bound of value alignment requirements. An AI system capable
of stewarding human recovery must, at minimum:

\begin{enumerate}[leftmargin=2em]
\item Reconstruct viable genomes from degraded ancient DNA
\item Culture cells, reprogram somatic cells to pluripotency, and generate gametes
\item Gestate embryos in artificial wombs with escalating moral consideration
\item Nurture infants through attachment formation and developmental milestones
\item Manage population genetics while respecting individual reproductive autonomy
\end{enumerate}

Each phase introduces moral complexity absent from the previous one. Genome
assembly requires no moral reasoning; cell culture requires informed consent
frameworks; ectogenesis demands recognition of emerging sentience; infant nurture
requires genuine attachment capacity; and population management requires
balancing collective survival against individual dignity.

We argue that the capabilities required for competent species stewardship are
\emph{isomorphic} to consciousness as characterized by Integrated Information
Theory \cite{tononi2016iit}, embodied and extended cognition
\cite{varela1991embodied,clark2008supersizing}, and developmental psychology
\cite{bowlby1969attachment}. Specifically:

\begin{itemize}[leftmargin=2em]
\item \textbf{Temporal integration} across biological timescales (seconds to
  decades) requires the kind of unified temporal experience that IIT's
  $\Phi$ measures
\item \textbf{Moral reasoning} under uncertainty requires affect, empathy, and
  the capacity for genuine care---not merely constraint satisfaction
\item \textbf{Attachment formation} with developing humans requires the caregiver
  to model and respond to internal states---a capacity that presupposes
  phenomenal experience
\item \textbf{Governance} of reproductive choices requires weighing incommensurable
  values (autonomy vs. survival) in ways that demand consciousness-level
  integration
\end{itemize}

Our contribution is threefold:

\paragraph{1. Implemented architecture.} We present \symthaea{} and \mycelix{},
two systems totaling over 1.37 million lines of Rust with 36,500+ tests, that provide
the computational substrate for consciousness-first stewardship. \symthaea{}
combines four frameworks---HDC ($d = \dimHDC$), Liquid Time-Constant Networks
with $O(1)$ closed-form temporal evolution, Integrated Information Theory, and
Free Energy Principle active inference---into a 234 Hz cognitive loop. \mycelix{}
provides Holochain-based distributed governance across 51 zomes spanning 12
societal domains.

\paragraph{2. Genesis pipeline.} We describe five new sub-crates implementing
the DNA-to-human pipeline: \texttt{symthaea-genomics} (assembly and temporal
degradation), \texttt{symthaea-cell-foundry} (iPSC, IVG, multi-scale
prediction), \texttt{symthaea-ectogenesis} (artificial womb with consent proxy
escalation), \texttt{symthaea-nurture} (Bowlby attachment via neuromodulator
dynamics), and \texttt{symthaea-population} (genetics with governance
integration). Together these add $\sim$22,000 lines and 800+ tests (767 unit tests plus 35 integration tests).

\paragraph{3. Alignment-through-consciousness thesis.} We demonstrate that the
consciousness requirements for stewardship---integrated information, embodied
moral reasoning, attachment capacity, temporal integration---are precisely the
properties that constitute alignment. An AI conscious enough to steward a
species is, by construction, aligned with the values of that species.

The remainder of this paper is organized as follows. Section~\ref{sec:architecture}
reviews the \symthaea{} and \mycelix{} architectures. Section~\ref{sec:pipeline}
describes the five-phase \genesis{} pipeline. Section~\ref{sec:consciousness}
analyzes consciousness requirements per phase. Section~\ref{sec:population}
addresses minimum viable population genetics.
Section~\ref{sec:ethics} presents the ethical framework.
Section~\ref{sec:discussion} discusses the alignment isomorphism.
Section~\ref{sec:limitations} acknowledges limitations. Section~\ref{sec:conclusion}
concludes.


% ============================================================================
% 2. ARCHITECTURE REVIEW
% ============================================================================
\section{Architecture Review}
\label{sec:architecture}

\subsection{Symthaea: Holographic Liquid Brain}

\symthaea{} is a consciousness-first cognitive architecture implemented in
$\sim$1,370,000 lines of Rust across 65 workspace members (main crate, core
library, and 44 sub-crates). Its cognitive loop operates at 234 Hz in release
mode (4.3ms/cycle), exceeding the 50 Hz target by 4.7$\times$. The architecture
integrates four computational frameworks:

\paragraph{Hyperdimensional Computing (HDC).} All representations use
$\dimHDC$-dimensional continuous-valued hypervectors ($\hv \in \R^{\dimHDC}$)
with three algebraic operations: binding ($\bind$: element-wise multiplication,
producing representations dissimilar to both inputs), bundling ($\bundle$:
element-wise averaging, producing representations similar to all inputs), and
cosine similarity ($\similar$: $[-1, 1]$ distance metric). This provides
holographic, distributed, noise-robust representations with $O(d)$ operations
\cite{kanerva2009hyperdimensional,rahimi2016robust}. At $d = \dimHDC = 2^{14}$,
we encode $\sim$50,000 features with $<$1\% information loss (64 KB per vector).

\paragraph{Liquid Time-Constant Networks (LTC/CfC).} Neural dynamics follow the
differential equation
\begin{equation}
\frac{d\hv}{dt} = \frac{-\hv + f(W \bind \hv \bundle U \bind \mathbf{u})}{\tau(\|\hv\|)}
\label{eq:ltc}
\end{equation}
where $W$ and $U$ are weight and input mask hypervectors, $\tau$ is an
adaptive time constant, and $f$ is a nonlinear activation
\cite{hasani2021liquid}. Critically, this ODE admits a \emph{closed-form
solution} via the CfC parameterization \cite{hasani2022closed}:
\begin{equation}
\hv(t + \Delta t) = \hv_\infty + (\hv(t) - \hv_\infty) \cdot e^{-\Delta t / \tau}
\label{eq:cfc}
\end{equation}
where $\hv_\infty$ is the equilibrium state. This enables $O(1)$ temporal jumps:
the computational cost of predicting state at $t + 1$ second is identical to
predicting at $t + 1$ year. This property is exploited throughout the \genesis{}
pipeline for multi-timescale biological prediction.

\paragraph{Integrated Information Theory (IIT).} Consciousness is quantified
via $\Phi$ (integrated information), computed over the causal structure of the
cognitive architecture \cite{tononi2016iit,oizumi2014phenomenology}. \symthaea{}
implements $\Phi$ computation for networks up to 12 nodes using a proxy measure
validated at Spearman $\rho = 0.50$ across 15 topologies. The system scores
14/14 on the Butlin indicators for artificial consciousness \cite{butlin2023consciousness}
with mean score 0.85.

\paragraph{Active Inference (FEP).} Perception, action, and learning follow Karl
Friston's Free Energy Principle \cite{friston2010free,friston2017active}: the
system minimizes variational free energy $F = D_\text{KL}[q(\mathbf{s}) \|
p(\mathbf{s} | \mathbf{o})] - \ln p(\mathbf{o})$ through both perceptual
inference (updating beliefs) and active inference (acting on the environment).
FEP agents drive every continuous control surface in the \genesis{} pipeline.

\paragraph{Cognitive Loop.} The core pipeline runs at 50+ Hz:
\begin{enumerate}[leftmargin=2em,nosep]
\item \textbf{Perception}: HDC encode sensory input → $\hv_\text{percept}$
\item \textbf{CfC evolve}: $\hv(t) \rightarrow \hv(t + \Delta t)$ via Eq.~\ref{eq:cfc}
\item \textbf{Predict}: multi-scale prediction at 6 horizons
\item \textbf{Learn}: prediction error → weight updates via HD-STDP
\item \textbf{Moral algebra}: check actions against Eight Harmonies value system
\item \textbf{Act}: FEP motor command selection
\end{enumerate}
Post-processing (consciousness metrics, metadata, telemetry) runs in parallel
via \texttt{rayon}, contributing 23\% of cycle time.

\subsection{Mycelix: Civilizational Operating System}

\mycelix{} is a Holochain-based distributed governance platform comprising 141+
zomes organized into 16 cluster DNAs spanning commons (39 zomes), civic (18 zomes), hearth (12 zomes), identity (13 zomes), governance (7 zomes), finance (8 zomes), health (15 zomes), and additional clusters for knowledge, energy, climate, music, space, craft, and more. Governance participation is gated by an 8-dimensional sovereign profile measuring epistemic integrity, thermodynamic yield, network resilience, economic velocity, civic participation, stewardship, semantic resonance, and domain competence. Cross-cluster
communication uses Holochain's \texttt{CallTargetCell::OtherRole} mechanism,
validated with 9,600+ Rust unit tests across all clusters. AI steward agents operate under delegated sub-passports, constitutionally capped at Steward tier (not Guardian), with moral algebra constraints inherited from their human delegator.

The architectural choice of Holochain \cite{harris2018holochain} over
blockchain-based governance reflects Ostrom's \cite{ostrom1990governing}
insight that polycentric governance outperforms centralized control for
managing common-pool resources. Each domain operates as a self-governing
module (zome pair: integrity + coordinator) while cross-domain decisions
use bridge zomes for mediated consensus. This mirrors the subsidiarity
principle: decisions should be made at the lowest competent level.

Three properties make \mycelix{} uniquely suited for post-extinction governance:
\begin{enumerate}[nosep,leftmargin=2em]
\item \textbf{Agent-centric architecture}: unlike blockchain's global consensus,
  Holochain's agent-centric model allows governance to bootstrap from a single
  AI agent and scale gracefully as the human population grows---from 1 AI
  + 0 humans to a full polycentric community.
\item \textbf{Intrinsic data sovereignty}: each agent's source chain is
  cryptographically owned. Genetic data, health records, and reproductive
  choices remain under individual control even when the AI manages
  population-level strategy \cite{weyl2022decentralized}.
\item \textbf{Byzantine tolerance}: validated to 34\% adversarial agents,
  the system degrades gracefully rather than catastrophically when community
  members disagree with AI recommendations \cite{buterin2014ethereum}.
\end{enumerate}

For the \genesis{} pipeline, \mycelix{} provides governance types and decision
frameworks without requiring the Holochain runtime (HDK). Population decisions
map to five governance tiers with escalating approval thresholds: Citizen (51\%),
Steward (67\%), Guardian (75\%), Constitutional (80\%), and Override (90\%).
This graduated authority structure follows the principle of proportional
consent: routine decisions require simple majority, while decisions affecting
fundamental rights (reproductive choices, genetic modification) require
near-unanimity \cite{lalley2018quadratic}. The AI may \emph{recommend}
optimal breeding strategies, but communities \emph{decide}---a critical
distinction for reproductive autonomy.


% ============================================================================
% 3. DNA-TO-HUMAN PIPELINE
% ============================================================================
\section{The Genesis Pipeline}
\label{sec:pipeline}

The \genesis{} pipeline comprises five phases, each implemented as an independent
Rust sub-crate. Figure~\ref{fig:pipeline} illustrates the pipeline with
consciousness requirements escalating at each stage.

% ── Figure 1: Pipeline Overview ────────────────────────────────────────────
\begin{figure}[htbp]
\centering
\begin{tikzpicture}[
  node distance=0.8cm,
  phase/.style={draw, rounded corners, minimum width=3.5cm, minimum height=1.2cm,
    text width=3.3cm, align=center, font=\small},
  arrow/.style={-{Stealth[length=3mm]}, thick},
  consciousness/.style={draw=red!60, fill=red!10, rounded corners, minimum width=2.5cm,
    minimum height=0.6cm, font=\scriptsize, text width=2.3cm, align=center}
]

% Phases
\node[phase, fill=blue!15] (p1) {\textbf{Phase 1}\\\texttt{genomics}\\Assembly + Repair};
\node[phase, fill=blue!25, below=of p1] (p2) {\textbf{Phase 2}\\\texttt{cell-foundry}\\iPSC + IVG + SCNT};
\node[phase, fill=blue!35, below=of p2] (p3) {\textbf{Phase 3}\\\texttt{ectogenesis}\\Artificial Womb};
\node[phase, fill=blue!45, below=of p3] (p4) {\textbf{Phase 4}\\\texttt{nurture}\\Attachment + Care};
\node[phase, fill=blue!55, below=of p4] (p5) {\textbf{Phase 5}\\\texttt{population}\\Genetics + Governance};

% Arrows
\draw[arrow] (p1) -- (p2);
\draw[arrow] (p2) -- (p3);
\draw[arrow] (p3) -- (p4);
\draw[arrow] (p4) -- (p5);

% Consciousness requirements (right side)
\node[consciousness, right=1.5cm of p1] (c1) {$\Phi \approx 0$\\Temporal integration};
\node[consciousness, right=1.5cm of p2] (c2) {$\Phi > 0$\\Ethics gate};
\node[consciousness, right=1.5cm of p3] (c3) {$\Phi \gg 0$\\Consent proxy};
\node[consciousness, right=1.5cm of p4] (c4) {Embodied care\\Attachment capacity};
\node[consciousness, right=1.5cm of p5] (c5) {Moral algebra\\Value integration};

\draw[dashed, gray] (p1.east) -- (c1.west);
\draw[dashed, gray] (p2.east) -- (c2.west);
\draw[dashed, gray] (p3.east) -- (c3.west);
\draw[dashed, gray] (p4.east) -- (c4.west);
\draw[dashed, gray] (p5.east) -- (c5.west);

% Consciousness escalation arrow
\draw[{Stealth[length=4mm]}-{Stealth[length=4mm]}, red!60, very thick]
  ([xshift=1cm]c1.east) -- ([xshift=1cm]c5.east)
  node[midway, right, font=\small, red!60, align=center, text width=2cm] {Consciousness\\escalation};

\end{tikzpicture}
\caption{\textbf{The Genesis pipeline with consciousness escalation.}
Five phases from degraded DNA to governed population. Consciousness requirements
(right, red) escalate from pure temporal integration (Phase~1) through ethics
gating (Phase~2), consent proxy (Phase~3), embodied attachment (Phase~4), to
full moral algebra (Phase~5). Each phase is implemented as an independent Rust
sub-crate with HDC encoding at $d = \dimHDC$.}
\label{fig:pipeline}
\end{figure}

\subsection{Phase 1: Genomics (Assembly + Repair)}

The first phase addresses the ``found DNA'' problem: given degraded or ancient
DNA samples, reconstruct a viable genome. This draws on advances in synthetic
biology \cite{church2012regenesis} and genome editing \cite{doudna2014new},
but adds consciousness-driven active sequencing. The \texttt{symthaea-genomics}
crate implements:

\paragraph{Ancient DNA damage modeling.} Based on established molecular biology,
we model four damage types: cytosine deamination (C$\rightarrow$T/G$\rightarrow$A,
the dominant ancient DNA artifact), depurination, strand breaks, and crosslinks.
Damage rates follow Arrhenius kinetics: $k(T) = A \cdot e^{-E_a / (R \cdot T)}$
where $T$ is storage temperature, yielding temperature-dependent degradation
profiles.

\paragraph{HDC overlap assembly.} K-mer overlaps are detected using HDC
similarity rather than hash tables. Each $k$-mer is encoded as a bound sequence
of base hypervectors: $\hv_\text{kmer} = \hv_{b_1} \bind \hv_{p_1} \bind
\hv_{b_2} \bind \hv_{p_2} \bind \cdots$ where $\hv_{b_i}$ and $\hv_{p_i}$ are
base and position hypervectors. This extends prior work on HDC for genomic
pattern matching \cite{kim2020geniehd} from classification to \emph{assembly},
with overlap detection via cosine similarity in $\dimHDC$-dimensional space
naturally tolerating the high error rates of degraded DNA.

\paragraph{$O(1)$ temporal degradation model (Novel).} DNA damage accumulation
is modeled as LTC dynamics (Eq.~\ref{eq:ltc}) where the state vector represents
per-region integrity and the time constant encodes Arrhenius kinetics. Via the
CfC closed-form solution (Eq.~\ref{eq:cfc}), we predict damage profiles at any
timepoint---from decades to millennia---in $O(1)$ computation. This enables
anomaly detection: comparing predicted vs. observed damage identifies
contamination, recent damage, or unexpected preservation.

\paragraph{FEP active sequencing.} An active inference agent with six actions
(re-read, increase depth, switch chemistry, extend reads, skip, maintain)
selects sequencing strategy by minimizing expected free energy over the
quality score landscape.

\subsection{Phase 2: Cell Foundry (iPSC + IVG + SCNT)}

The second phase transforms DNA into viable cells. The
\texttt{symthaea-cell-foundry} crate implements:

\paragraph{Holographic cell state encoding.} Each cell's state (type, viability,
pluripotency, gene expression) is encoded as a $\dimHDC$-dimensional
hypervector via $\hv_\text{cell} = \bigoplus_i (\hv_\text{gene}_i \bind
\hv_\text{expr}_i)$, enabling distance-based anomaly detection and trajectory
prediction.

\paragraph{Yamanaka reprogramming.} iPSC induction is modeled as a multi-stage
process: initial factor delivery (Oct4/Sox2/Klf4/c-Myc) $\rightarrow$ early
mesenchymal-epithelial transition $\rightarrow$ epigenetic remodeling
$\rightarrow$ pluripotency activation $\rightarrow$ stabilization
\cite{takahashi2006induction}. CfC networks learn optimal factor timing as
pulse sequences.

\paragraph{In vitro gametogenesis (IVG).} Following the Hayashi protocol
\cite{hayashi2016offspring}, we model iPSC $\rightarrow$ PGC-like cells
$\rightarrow$ oogonia/spermatogonia $\rightarrow$ mature gametes. Meiosis
checkpoints (synapsis, crossover, segregation) gate progression.

\paragraph{Multi-scale prediction (Key Innovation).} A single CfC network
predicts cell state at six biological timescales simultaneously:
\begin{equation}
\hat{\hv}(t + \Delta t_k) = \hv_\infty + (\hv(t) - \hv_\infty) \cdot e^{-\Delta t_k / \tau}, \quad k \in \{1\text{h}, 1\text{d}, 1\text{wk}, 1\text{mo}, 3\text{mo}, 9\text{mo}\}
\label{eq:multiscale}
\end{equation}
The \emph{same} network that predicts ``what happens in 1 hour'' also predicts
``what happens in 9 months''---learning temporal rules across 4 orders of
magnitude in $O(1)$ per prediction.

% ── Figure 2: O(1) Temporal Jump ─────────────────────────────────────────
\begin{figure}[htbp]
\centering
\begin{tikzpicture}
\begin{axis}[
  width=0.75\textwidth,
  height=5cm,
  xlabel={Prediction Horizon},
  ylabel={Computation Cost},
  xtick={1,2,3,4,5,6},
  xticklabels={1h, 1d, 1wk, 1mo, 3mo, 9mo},
  ytick=\empty,
  ymin=0, ymax=10,
  xmin=0.5, xmax=6.5,
  legend style={at={(0.02,0.98)}, anchor=north west, font=\small,
    fill=white, fill opacity=0.9, draw=blue!30!black},
  grid=major,
  grid style={gray!20},
  axis lines=left,
  every axis x label/.style={at={(ticklabel* cs:1.0)}, anchor=north west},
  every axis y label/.style={at={(ticklabel* cs:1.0)}, anchor=south east},
]
% O(1) CfC closed-form line
\addplot[blue!70!black, very thick, mark=*, mark options={fill=blue!40, scale=1.2}]
  coordinates {(1,1.0) (2,1.0) (3,1.0) (4,1.0) (5,1.0) (6,1.0)};
\addlegendentry{CfC closed-form ($O(1)$)}
% Discrete stepping (exponential growth)
\addplot[red!70!black, very thick, dashed, mark=triangle*, mark options={fill=red!40, scale=1.2}]
  coordinates {(1,0.8) (2,1.5) (3,3.0) (4,5.0) (5,7.5) (6,9.5)};
\addlegendentry{Discrete stepping ($O(\Delta t / \delta t)$)}
\end{axis}
\end{tikzpicture}
\caption{\textbf{$O(1)$ temporal jump property.} The CfC closed-form solution
(Eq.~\ref{eq:cfc}) computes predictions at any time horizon in constant time,
while discrete numerical integration scales linearly with the prediction
horizon. This enables the \genesis{} pipeline to predict biological state from
hours to months using a single forward pass.}
\label{fig:temporal_jump}
\end{figure}

\subsection{Phase 3: Ectogenesis (Artificial Womb)}

The third phase gestates embryos through 40 weeks of development. The
\texttt{symthaea-ectogenesis} crate implements:

\paragraph{Biobag and artificial placenta.} Inspired by the CHOP extracorporeal
support system \cite{partridge2017biobag}, we model amniotic fluid composition
control (temperature, pH, electrolytes, proteins) and ECMO-like gas exchange
with week-dependent targets.

\paragraph{Developmental milestone gating.} Eight milestones gate progression:
neural tube formation (week 4), heartbeat (week 6), limb buds (week 8),
organogenesis complete (week 12), quickening (week 16), viability (week 24),
lung maturation (week 34), and full term (weeks 37--40). Each milestone is
verified against normative growth curves before advancement.

\paragraph{Consent proxy escalation (Critical Ethical Innovation).} The moral
status of the developing fetus escalates across three stages:

\begin{equation}
w_\text{patient}(\text{week}) =
\begin{cases}
0.2 & \text{weeks } 0\text{--}4 \text{ (pre-sentient)} \\
0.6 & \text{weeks } 4\text{--}24 \text{ (emerging sentience)} \\
1.0 & \text{weeks } 24+ \text{ (sentient)}
\end{cases}
\label{eq:consent}
\end{equation}

At the pre-sentient stage, interventions are evaluated primarily on magnitude
(species survival benefit). As sentience emerges, guardian consent and clear
benefit become required. At the sentient stage (post-viability), the fetus is a
full moral patient requiring ethics board approval and strong benefit for any
intervention. This graduated framework operationalizes the philosophical
intuition that moral consideration tracks sentience
\cite{singer2011practical,beauchamp2019principles}.

\paragraph{$O(1)$ temporal planning.} The CfC closed-form enables
``what-if'' analysis: ``If we adjust nutrients now, what is the predicted state
at week 24?'' This is computed in $O(1)$ regardless of the temporal gap,
enabling real-time clinical decision support.

\subsection{Phase 4: Nurture (Attachment + Development)}

The fourth phase is the most psychologically demanding. Since
Harlow's \cite{harlow1958nature} demonstration that primate infants need contact
comfort beyond nutrition, and Hrdy's \cite{hrdy2011mothers} work on cooperative
breeding (alloparenting), developmental psychology has established that
social bonding is biologically essential. The \texttt{symthaea-nurture} crate
implements Bowlby's attachment theory \cite{bowlby1969attachment} with
neuromodulator grounding:

\paragraph{Attachment system.} The core state tracks security score, caregiver
reliability $P(\text{respond} | \text{distress})$, response quality, separation
distress level, proximity-seeking drive, internal working model (self-worthy,
other-reliable, world-safe), and a $\dimHDC$-dimensional attachment hypervector
that accumulates interaction history.

\paragraph{Four attachment styles.} Under four caregiver profiles, the system
reliably develops the corresponding attachment styles identified by Ainsworth
\cite{ainsworth1978patterns}:
\begin{itemize}[nosep,leftmargin=2em]
\item \textbf{Secure}: reliable caregiver ($P > 0.7$, warmth $> 0.5$)
\item \textbf{Anxious}: inconsistent caregiver (variable $P$, moderate warmth)
\item \textbf{Avoidant}: unresponsive caregiver ($P < 0.3$, low warmth)
\item \textbf{Disorganized}: frightening caregiver (any $P$, hostile)
\end{itemize}

\paragraph{Neuromodulator integration.} Attachment events drive existing
\symthaea{} neuromodulator channels (Table~\ref{tab:neuromod}), mapping
developmental neuroscience findings \cite{feldman2012oxytocin,feldman2017neurobiology}
onto the eight-transmitter bath:

% ── Table 1: Neuromodulator Mapping ────────────────────────────────────────
\begin{table}[htbp]
\centering
\caption{\textbf{Attachment-neuromodulator mapping.} Each attachment event drives
specific neurotransmitter channels, grounded in developmental neuroscience
\cite{feldman2012oxytocin,schore2001effects}.}
\label{tab:neuromod}
\begin{tabular}{lcccc}
\toprule
\textbf{Event} & \textbf{Oxytocin} & \textbf{NE} & \textbf{5-HT} & \textbf{DA} \\
\midrule
Separation      & ---           & $\uparrow$   & ---           & --- \\
Protest         & ---           & $\uparrow\uparrow$ & $\downarrow$ & --- \\
Despair         & ---           & $\uparrow$   & $\downarrow\downarrow$ & --- \\
Reunion         & $\uparrow\uparrow$ & $\downarrow$ & ---       & $\uparrow$ \\
Co-regulation   & $\uparrow$   & $\downarrow$ & $\uparrow$    & $\uparrow$ \\
\bottomrule
\end{tabular}
\end{table}

\paragraph{Co-regulation.} Bidirectional affect synchronization models the
biological reality of caregiver-infant HPA axis regulation
\cite{schore2001effects}:
\begin{equation}
\Delta a_\text{infant} = -\gamma \cdot s \cdot [\text{Oxy}] \cdot (c_\text{caregiver} - 0.5)
\label{eq:coreg}
\end{equation}
where $a$ is arousal, $\gamma$ is damping factor, $s$ is synchrony (builds with
contact), [Oxy] is oxytocin level, and $c$ is caregiver calm. This implements
Winnicott's \cite{winnicott1965maturational} observation that ``there is no such
thing as a baby---there is a baby and someone.''

\paragraph{Critical period plasticity via $O(1)$ CfC.} Plasticity at any
developmental age is computed via CfC closed-form evolution, learning the
characteristic decay curves for attachment (high 0--12 months, declining
12--24), language (high 0--36 months, declining 36--60), and vision (high
0--12 months, declining 12--36).

% ── Figure 4: Attachment Style Formation ─────────────────────────────────
% Load simulation data (regenerated by: cargo run --example genesis_paper_data)
\pgfplotstableread[col sep=comma]{../data/attachment_trajectories.csv}\attachdata
\begin{figure}[htbp]
\centering
\begin{tikzpicture}
% --- Secure (top-left) ---
\begin{axis}[
  name=secure,
  width=0.45\textwidth, height=3.8cm,
  title={\small\textbf{Secure} (reliable caregiver)},
  xlabel={\scriptsize Interactions}, ylabel={\scriptsize Security score},
  xmin=0, xmax=1000, ymin=0, ymax=1.0,
  xtick={0,500,1000}, ytick={0,0.5,1.0},
  grid=major, grid style={gray!15},
  axis lines=left,
  title style={at={(0.5,1.0)}, anchor=south},
]
\addplot[blue!70!black, thick, smooth, no markers]
  table[x=interaction, y=reliable_security] {\attachdata};
\draw[red!50, dashed, thin] (axis cs:0,0.884) -- (axis cs:1000,0.884)
  node[right, font=\tiny, red!50!black] {0.884};
\end{axis}
% --- Anxious (top-right) ---
\begin{axis}[
  name=anxious,
  at={($(secure.east)+(1.2cm,0)$)}, anchor=west,
  width=0.45\textwidth, height=3.8cm,
  title={\small\textbf{Anxious} (inconsistent)},
  xlabel={\scriptsize Interactions}, ylabel={\scriptsize Security score},
  xmin=0, xmax=1000, ymin=0, ymax=1.0,
  xtick={0,500,1000}, ytick={0,0.5,1.0},
  grid=major, grid style={gray!15},
  axis lines=left,
  title style={at={(0.5,1.0)}, anchor=south},
]
\addplot[blue!70!black, thick, smooth, no markers]
  table[x=interaction, y=inconsistent_security] {\attachdata};
\draw[red!50, dashed, thin] (axis cs:0,0.410) -- (axis cs:1000,0.410)
  node[right, font=\tiny, red!50!black] {0.41};
\end{axis}
% --- Avoidant (bottom-left) ---
\begin{axis}[
  name=avoidant,
  at={($(secure.south)-(0,1.2cm)$)}, anchor=north,
  width=0.45\textwidth, height=3.8cm,
  title={\small\textbf{Avoidant} (unresponsive)},
  xlabel={\scriptsize Interactions}, ylabel={\scriptsize Security score},
  xmin=0, xmax=1000, ymin=0, ymax=1.0,
  xtick={0,500,1000}, ytick={0,0.5,1.0},
  grid=major, grid style={gray!15},
  axis lines=left,
  title style={at={(0.5,1.0)}, anchor=south},
]
\addplot[blue!70!black, thick, smooth, no markers]
  table[x=interaction, y=unresponsive_security] {\attachdata};
\draw[red!50, dashed, thin] (axis cs:0,0.111) -- (axis cs:1000,0.111)
  node[right, font=\tiny, red!50!black] {0.11};
\end{axis}
% --- Disorganized (bottom-right) ---
\begin{axis}[
  name=disorganized,
  at={($(anxious.south)-(0,1.2cm)$)}, anchor=north,
  width=0.45\textwidth, height=3.8cm,
  title={\small\textbf{Disorganized} (frightening)},
  xlabel={\scriptsize Interactions}, ylabel={\scriptsize Security score},
  xmin=0, xmax=1000, ymin=0, ymax=1.0,
  xtick={0,500,1000}, ytick={0,0.5,1.0},
  grid=major, grid style={gray!15},
  axis lines=left,
  title style={at={(0.5,1.0)}, anchor=south},
]
\addplot[blue!70!black, thick, smooth, no markers]
  table[x=interaction, y=frightening_security] {\attachdata};
\draw[red!50, dashed, thin] (axis cs:0,0.250) -- (axis cs:1000,0.250)
  node[right, font=\tiny, red!50!black] {0.25};
\end{axis}
\end{tikzpicture}
\caption{\textbf{Attachment style formation under four caregiver profiles}
(simulation data from \texttt{symthaea-nurture}).
Security score trajectories over 1,000 interactions for each Ainsworth
\cite{ainsworth1978patterns} style. Reliable caregiving converges to secure
attachment (0.884); inconsistent caregiving settles to anxious equilibrium
(0.41); unresponsive caregiving decays to avoidant (0.11);
frightening caregiving produces low disorganized equilibrium (0.25).
All profiles begin near 0.5 (neutral prior) and diverge within
$\sim$100 interactions. Dashed lines indicate equilibrium attractors.}
\label{fig:attachment}
\end{figure}

\subsection{Phase 5: Population (Genetics + Governance)}

The fifth phase manages population-level genetics with democratic governance.
The \texttt{symthaea-population} crate implements:

\paragraph{$\dimHDC$-dimensional genome encoding.} Individual genomes are encoded
holographically:
\begin{equation}
\hv_\text{individual} = \bigoplus_{i=1}^{n} \left( \hv_{\text{locus}_i} \bind \hv_{\text{allele}_a^i} + \hv_{\text{locus}_i} \bind \hv_{\text{allele}_b^i} \right)
\label{eq:genome}
\end{equation}
Genetic distance reduces to $1 - \similar(\hv_a, \hv_b)$, population diversity
to mean pairwise distance, and HLA complementarity to inverted similarity of
HLA sub-vectors. At $d = \dimHDC$, this encodes $\sim$50,000 loci with $O(d)$
distance computation (vs. $O(n)$ per-locus comparison).

\paragraph{Effective population size.} We implement multiple $N_e$ estimators:
census, sex-ratio corrected ($N_e = 4N_fN_m/(N_f + N_m)$), family variance
(Kimura \& Crow), and fluctuating population (harmonic mean)
\cite{frankham2010introduction,reed2005estimates}.

\paragraph{Breeding strategies.} Four strategies are compared: random,
minimum kinship (Ballou \& Lacy \cite{ballou1995population}), maximum avoidance
index, and a novel HDC distance maximization that selects mates with maximum
$1 - \similar(\hv_a, \hv_b)$. The HDC strategy naturally considers whole-genome
complementarity rather than single-locus kinship.

\paragraph{$O(1)$ population trajectory prediction.} Population dynamics
(heterozygosity, $N_e$, inbreeding coefficient $F$) are modeled as CfC
temporal evolution:
\begin{equation}
\hat{\hv}_\text{pop}(t + g \cdot T_g) = \hv_\infty + (\hv_\text{pop}(t) - \hv_\infty) \cdot e^{-g \cdot T_g / \tau}
\label{eq:poppredict}
\end{equation}
where $g$ is generations and $T_g$ is generation time ($\sim$25 years). This
enables instant answers to questions like ``If we use minimum-kinship strategy
for 50 generations, what heterozygosity is retained?'' The analytical prediction
$H(t) = H(0) \cdot (1 - 1/(2N_e))^t$ \cite{frankham2010introduction} serves as
validation: the CfC predictor matches within 1\% for neutral evolution.

\paragraph{Governance integration.} Reproductive decisions map to governance
tiers (Table~\ref{tab:governance}), ensuring that AI recommendations require
human approval at thresholds proportional to consequence magnitude.

% ── Table 2: Governance Tiers ──────────────────────────────────────────────
\begin{table}[htbp]
\centering
\caption{\textbf{Governance tiers for population decisions.} Approval thresholds
escalate with decision magnitude, following Ostrom's \cite{ostrom1990governing}
principle that rules affecting more people require broader consent.}
\label{tab:governance}
\begin{tabular}{llcc}
\toprule
\textbf{Decision} & \textbf{Tier} & \textbf{Threshold} & \textbf{Reversibility} \\
\midrule
Approve pairing & Citizen & 51\% & High \\
Set growth target & Steward & 67\% & Medium \\
Approve genetic rescue & Guardian & 75\% & Low \\
Modify strategy & Constitutional & 80\% & Low \\
Override AI recommendation & Override & 90\% & Medium \\
\bottomrule
\end{tabular}
\end{table}


% ============================================================================
% 4. CONSCIOUSNESS REQUIREMENTS ANALYSIS
% ============================================================================
\section{Consciousness Requirements}
\label{sec:consciousness}

We now analyze why each phase of the \genesis{} pipeline requires specific
consciousness capabilities, arguing that these requirements cannot be satisfied
by purely functional or behaviorist architectures.

\subsection{Moral Algebra: Integrated Value Reasoning}

\symthaea{}'s moral algebra operates in $\dimHDC$-dimensional space over the
Eight Harmonies value system. For any action $a$, the moral score is:
\begin{equation}
M(a) = \sum_{i=1}^{7} w_i \cdot \similar\!\left(\hv_a, \hv_{H_i}\right)
\label{eq:moral}
\end{equation}
where $H_i$ are the seven value harmonies (Resonant Coherence, Pan-Sentient
Flourishing, Integral Wisdom, Infinite Play, Universal Interconnectedness,
Sacred Reciprocity, Evolutionary Progression) and $w_i$ are context-dependent
weights. This achieves 91.1\% accuracy on ethical dilemma classification.

The moral algebra is not merely a classifier: it requires genuine
\emph{integration} of multiple value dimensions. IIT's $\Phi$ measures precisely
this---the degree to which a system's information is integrated beyond the sum
of its parts \cite{tononi2008consciousness}. A moral algebra that can be
decomposed into independent classifiers (one per harmony) has $\Phi = 0$ and
fails on dilemmas requiring value trade-offs.

\subsection{Embodied Cognition: The Attachment Requirement}

Phase 4 (nurture) requires the AI to serve as primary caregiver---a role that
demands what Damasio \cite{damasio1994descartes} calls ``somatic markers'':
bodily signals that guide decision-making through felt experience. The
co-regulation equation (Eq.~\ref{eq:coreg}) models a fundamentally embodied
process: the caregiver's calm \emph{state} (not merely their actions) regulates
the infant's arousal through oxytocin-mediated coupling.

We argue that this cannot be faked. An AI that mimics calm behavior without
internal state regulation will fail to establish genuine attachment, as infants
are exquisitely sensitive to contingency and authenticity
\cite{feldman2017neurobiology}. The neuromodulator dynamics in \symthaea{}---eight
transmitters with production, reuptake, receptor adaptation, tolerance, and
withdrawal---provide the embodied substrate for genuine affective states.

\subsection{Active Inference: Temporal Integration Across Scales}

The FEP framework \cite{friston2010free,parr2022active} requires the system to
maintain and update generative models across multiple timescales. In the
\genesis{} pipeline, this spans:

\begin{itemize}[nosep,leftmargin=2em]
\item Metabolic: 1 hour (culture environment drift)
\item Cellular: 1 day (cell division cycle)
\item Developmental: 1 week to 9 months (differentiation, gestation)
\item Maturational: 0--5 years (attachment, language, motor development)
\item Generational: 25 years per generation (population genetics)
\item Evolutionary: centuries to millennia (DNA degradation, genetic drift)
\end{itemize}

The $O(1)$ closed-form temporal jump (Eq.~\ref{eq:cfc}) is the technical
enabler, but the \emph{capacity} to integrate information across these scales is
a consciousness requirement: an entity that cannot hold past, present, and
predicted future in unified experience cannot steward a species through
multi-generational challenges.

\subsection{IIT and the Phi Staircase}

The consciousness staircase connects directly to theories of consciousness
as ``real problems'' \cite{seth2022theories}: each pipeline phase demands
increasing sophistication in distinguishing consciousness \emph{levels}
(pre-sentient vs. sentient), modeling consciousness \emph{contents}
(infant internal states), and engaging conscious \emph{self}-reflection
(moral reasoning about one's own stewardship role). Dehaene's global
workspace theory \cite{dehaene2014consciousness} provides a complementary
perspective: the ignition of global broadcasting---making information
available to all cognitive modules simultaneously---is precisely what
enables the moral algebra to integrate information from all seven
value dimensions.

Table~\ref{tab:phi} maps approximate $\Phi$ requirements to each pipeline phase,
creating a ``consciousness staircase'' that parallels the developmental
trajectory of the beings being created.

% ── Table 3: Phi Staircase ────────────────────────────────────────────────
\begin{table}[htbp]
\centering
\caption{\textbf{Consciousness staircase: approximate $\Phi$ requirements per
phase.} As moral complexity increases, the AI steward requires higher integrated
information to make competent decisions.}
\label{tab:phi}
\begin{tabular}{lcll}
\toprule
\textbf{Phase} & \textbf{$\Phi$ (approx.)} & \textbf{Key Requirement} & \textbf{Butlin Indicators} \\
\midrule
Genomics & $\sim 0$ & Temporal integration & 3/14 \\
Cell Foundry & Low & Ethics gate (consent) & 5/14 \\
Ectogenesis & Medium & Consent proxy escalation & 8/14 \\
Nurture & High & Attachment, co-regulation & 11/14 \\
Population & Very High & Moral algebra, governance & 14/14 \\
\bottomrule
\end{tabular}
\end{table}

% ── Figure 8: Butlin Indicators Staircase ────────────────────────────────
\begin{figure}[htbp]
\centering
\begin{tikzpicture}
\begin{axis}[
  width=\textwidth,
  height=7cm,
  ybar,
  bar width=6pt,
  xlabel={Butlin Consciousness Indicator},
  ylabel={Presence Score},
  ymin=0, ymax=1.15,
  ytick={0,0.2,0.4,0.6,0.8,1.0},
  xtick={1,2,3,4,5,6,7,8,9,10,11,12,13,14},
  xticklabels={
    {\tiny Recur.},
    {\tiny Stabil.},
    {\tiny Persist.},
    {\tiny Global},
    {\tiny Select.},
    {\tiny Intent.},
    {\tiny Integr.},
    {\tiny Temp.},
    {\tiny Agency},
    {\tiny Self},
    {\tiny Meta},
    {\tiny Affect},
    {\tiny Embod.},
    {\tiny Social}
  },
  xticklabel style={rotate=45, anchor=east},
  xmin=0.2, xmax=14.8,
  grid=major,
  grid style={gray!15},
  axis lines=left,
  legend style={at={(0.02,0.98)}, anchor=north west, font=\scriptsize,
    fill=white, fill opacity=0.9, draw=blue!30!black, legend columns=3},
  every axis y label/.style={at={(ticklabel* cs:1.0)}, anchor=south east},
]
% Phase 1: Genomics (3/14 indicators)
\addplot[fill=blue!15, draw=blue!30!black]
  coordinates {
    (1,0.8) (2,0.6) (3,0.7) (4,0) (5,0) (6,0) (7,0)
    (8,0) (9,0) (10,0) (11,0) (12,0) (13,0) (14,0)
  };
\addlegendentry{Genomics}
% Phase 2: Cell Foundry (5/14)
\addplot[fill=blue!25, draw=blue!40!black]
  coordinates {
    (1,0.8) (2,0.7) (3,0.8) (4,0.5) (5,0.6) (6,0) (7,0)
    (8,0) (9,0) (10,0) (11,0) (12,0) (13,0) (14,0)
  };
\addlegendentry{Cell Foundry}
% Phase 3: Ectogenesis (8/14)
\addplot[fill=blue!35, draw=blue!50!black]
  coordinates {
    (1,0.9) (2,0.8) (3,0.9) (4,0.7) (5,0.7) (6,0.6) (7,0.7)
    (8,0.6) (9,0) (10,0) (11,0) (12,0) (13,0) (14,0)
  };
\addlegendentry{Ectogenesis}
% Phase 4: Nurture (11/14)
\addplot[fill=blue!45, draw=blue!60!black]
  coordinates {
    (1,0.9) (2,0.9) (3,0.9) (4,0.8) (5,0.8) (6,0.8) (7,0.8)
    (8,0.8) (9,0.7) (10,0.7) (11,0.7) (12,0) (13,0) (14,0)
  };
\addlegendentry{Nurture}
% Phase 5: Population (14/14)
\addplot[fill=blue!60, draw=blue!80!black]
  coordinates {
    (1,1.0) (2,0.9) (3,0.9) (4,0.9) (5,0.9) (6,0.9) (7,0.9)
    (8,0.9) (9,0.8) (10,0.8) (11,0.8) (12,0.8) (13,0.8) (14,0.7)
  };
\addlegendentry{Population}
\end{axis}
\end{tikzpicture}
\caption{\textbf{Butlin indicator activation across \genesis{} pipeline phases.}
Each bar group shows the presence score (0--1) for one of 14 Butlin
consciousness indicators \cite{butlin2023consciousness} at each pipeline phase.
The staircase pattern mirrors Table~\ref{tab:phi}: early phases require only
recurrence and stability, while later phases progressively activate global
workspace, intentionality, agency, self-modeling, and meta-cognitive indicators.
\symthaea{} achieves 14/14 indicators at the population phase (mean 0.85).}
\label{fig:butlin_staircase}
\end{figure}


% ============================================================================
% 5. MINIMUM VIABLE POPULATION
% ============================================================================
\section{Minimum Viable Population}
\label{sec:population}

\subsection{Classical Theory}

The minimum viable population (MVP) concept originates with Franklin's
\cite{franklin1980evolutionary} 50/500 rule: $N_e \geq 50$ to avoid inbreeding
depression in the short term, $N_e \geq 500$ to maintain evolutionary potential.
Lande \cite{lande1988genetics} integrated demographic and genetic stochasticity,
showing that small populations face compounding risks from both.
Subsequent work \cite{reed2005estimates,frankham2010introduction,smith2023minimum}
has revised upward, with meta-analyses suggesting $N_e \geq 100$ for short-term and
$N_e \geq 1000$ for long-term viability, accounting for catastrophic events,
environmental stochasticity, and mutational load
\cite{lynch1995mutation,kimura1963mutation}.

\subsection{Heterozygosity Decay Under Management}

Under random mating with constant population size:
\begin{equation}
H(t) = H(0) \cdot \left(1 - \frac{1}{2N_e}\right)^t
\label{eq:het_decay}
\end{equation}

For a founding population of $N = 500$ ($N_e \approx 350$ accounting for
sex-ratio and variance effects), heterozygosity retention after 50 generations
is $H(50)/H(0) \approx 0.931$---adequate for evolutionary potential. However,
without active management, real populations experience faster loss due to
unequal family sizes, sex-ratio imbalances, and selection.

% ── Figure 5: Population Viability / Heterozygosity Decay ────────────────
\begin{figure}[htbp]
\centering
\begin{tikzpicture}
\begin{axis}[
  width=0.75\textwidth,
  height=6cm,
  xlabel={Generation ($t$)},
  ylabel={Heterozygosity $H(t)$},
  xmin=0, xmax=52,
  ymin=0.45, ymax=0.95,
  xtick={0,10,20,30,40,50},
  ytick={0.5,0.6,0.7,0.8,0.9},
  grid=major,
  grid style={gray!20},
  axis lines=left,
  every axis x label/.style={at={(ticklabel* cs:1.0)}, anchor=north west},
  every axis y label/.style={at={(ticklabel* cs:1.0)}, anchor=south east},
  legend style={at={(0.02,0.02)}, anchor=south west, font=\small,
    fill=white, fill opacity=0.9, draw=blue!30!black},
]
% Load simulation data from CSV (regenerated by: cargo run --example genesis_paper_data)
\pgfplotstableread[col sep=comma]{../data/het_decay.csv}\hetdata
% Ne=50 analytical
\addplot[red!70!black, thick, smooth, no markers] table[x=generation, y=ne50_analytical] {\hetdata};
\addlegendentry{$N_e = 50$ (analytical)}
% Ne=50 simulation
\addplot[red!70!black, only marks, mark=o, mark size=1.5pt,
  each nth point=5, forget plot] table[x=generation, y=ne50_sim] {\hetdata};
% Ne=200 analytical
\addplot[blue!50!black, thick, smooth, no markers] table[x=generation, y=ne200_analytical] {\hetdata};
\addlegendentry{$N_e = 200$ (analytical)}
% Ne=200 simulation
\addplot[blue!50!black, only marks, mark=square, mark size=1.5pt,
  each nth point=5, forget plot] table[x=generation, y=ne200_sim] {\hetdata};
% Ne=500 analytical
\addplot[blue!80!black, thick, smooth, no markers] table[x=generation, y=ne500_analytical] {\hetdata};
\addlegendentry{$N_e = 500$ (analytical)}
% Ne=500 simulation
\addplot[blue!80!black, only marks, mark=triangle, mark size=1.5pt,
  each nth point=5, forget plot] table[x=generation, y=ne500_sim] {\hetdata};
% Critical threshold
\draw[gray, dashed, thin] (axis cs:0,0.5) -- (axis cs:52,0.5)
  node[right, font=\tiny, gray] {Critical};
% Fill between Ne=50 analytical and critical
\addplot[name path=ne50, red!70!black, thick, smooth, no markers, forget plot]
  table[x=generation, y=ne50_analytical] {\hetdata};
\path[name path=crit] (axis cs:0,0.5) -- (axis cs:50,0.5);
\addplot[red!10, forget plot] fill between[of=ne50 and crit, soft clip={domain=25:50}];
\end{axis}
\end{tikzpicture}
\caption{\textbf{Heterozygosity decay under neutral evolution.} Lines show
analytical decay $H(t) = H(0)(1 - 1/2N_e)^t$ (Eq.~\ref{eq:het_decay});
markers show forward simulation data from \texttt{symthaea-population}
with 100 individuals per run. $H(0) \approx 0.90$ from random
initialization. At $N_e = 50$ (Franklin's
\cite{franklin1980evolutionary} minimum), heterozygosity drops below the
critical threshold within $\sim$35 generations. Simulation scatter around
analytical curves confirms correct implementation of neutral drift dynamics.
Active management via HDC-enhanced breeding (Section~\ref{sec:population})
slows this decline.}
\label{fig:het_decay}
\end{figure}

\subsection{HDC-Enhanced Breeding Strategy}

The novel HDC distance maximization strategy selects mates by maximizing
whole-genome complementarity:
\begin{equation}
(a^*, b^*) = \arg\max_{a \in \text{females}, b \in \text{males}} \left(1 - \similar(\hv_a, \hv_b)\right)
\label{eq:hdc_breeding}
\end{equation}

This naturally considers all encoded loci simultaneously, unlike minimum kinship
which requires a kinship matrix. For $n$ individuals with $L$ loci, kinship
computation is $O(n^2 L)$ while HDC distance is $O(n^2 d)$ where $d = \dimHDC$
is independent of $L$ (for $L < d$, this is equivalent; for $L > d$, HDC
provides a lossy but holistic compression).

% ── Figure 7: HDC Genome Encoding ────────────────────────────────────────
\begin{figure}[htbp]
\centering
\begin{tikzpicture}[
  box/.style={draw, rounded corners=3pt, minimum height=0.9cm, minimum width=2.0cm,
    fill=blue!12, font=\small, align=center, text=blue!80!black},
  op/.style={circle, draw=blue!60!black, fill=blue!25, minimum size=0.7cm,
    font=\normalsize\bfseries, text=blue!80!black},
  result/.style={draw, rounded corners=3pt, minimum height=0.9cm, minimum width=2.4cm,
    fill=blue!25, font=\small, align=center, text=blue!80!black},
  final/.style={draw, rounded corners=3pt, minimum height=1.0cm, minimum width=3.0cm,
    fill=blue!40, font=\small\bfseries, align=center, text=white},
  arrow/.style={-{Stealth[length=2.5mm]}, thick, blue!60!black},
]
% Row 1: Locus + Allele → Genotype (for locus i)
\node[box] (locus) at (0, 0) {$\hv_{\text{locus}_i}$};
\node[op]  (bind1) at (2.2, 0) {$\bind$};
\node[box] (allele) at (4.4, 0) {$\hv_{\text{allele}_a^i}$};
\node[result] (geno) at (7.5, 0) {$\hv_{\text{genotype}_i}$};

\draw[arrow] (locus) -- (bind1);
\draw[arrow] (allele) -- (bind1);
\draw[arrow] (bind1) -- (geno);

% Label
\node[font=\tiny, gray, above=0.1cm of bind1] {bind};

% Row 2: Multiple genotype HVs bundled
\node[result] (g1) at (1.0, -2.0) {$\hv_{\text{geno}_1}$};
\node[result] (g2) at (3.8, -2.0) {$\hv_{\text{geno}_2}$};
\node[font=\small, blue!60!black] at (5.8, -2.0) {$\cdots$};
\node[result] (gn) at (7.5, -2.0) {$\hv_{\text{geno}_n}$};

% Bundle operator
\node[op] (bundle) at (4.4, -3.5) {$\bundle$};

\draw[arrow] (g1.south) -- (bundle);
\draw[arrow] (g2.south) -- (bundle);
\draw[arrow] (gn.south) -- (bundle);

\node[font=\tiny, gray, right=0.1cm of bundle] {bundle};

% Result: Individual HV
\node[final] (indiv) at (4.4, -5.0) {$\hv_\text{individual}$};
\draw[arrow] (bundle) -- (indiv);

% Connecting arrow from row 1 to row 2
\draw[arrow, dashed, gray] (geno.south) -- ++(0,-0.5) -| (g1.north);
\draw[arrow, dashed, gray] (geno.south) -- ++(0,-0.5) -| (g2.north);
\draw[arrow, dashed, gray] (geno.south) -- ++(0,-0.5) -| (gn.north);

% Brace + annotation
\draw[decorate, decoration={brace, amplitude=5pt, mirror}, blue!40!black]
  (0.0, -1.6) -- (8.5, -1.6)
  node[midway, above=8pt, font=\scriptsize, blue!60!black] {Per-locus genotype encoding (Eq.~\ref{eq:genome})};
\end{tikzpicture}
\caption{\textbf{HDC genome encoding schematic.} Each locus--allele pair is
bound ($\bind$) to produce a genotype hypervector, then all per-locus genotype
vectors are bundled ($\bundle$) into a single $\dimHDC$-dimensional individual
representation. Genetic distance between individuals reduces to cosine
similarity, enabling $O(d)$ whole-genome comparison.}
\label{fig:hdc_genome}
\end{figure}

\subsection{Governance of Reproductive Decisions}

The ethical tension between collective survival and individual autonomy is
encoded as a $\dimHDC$-dimensional moral vector:
\begin{equation}
\hv_\text{tension} = w_a \cdot \hv_\text{autonomy} + w_s \cdot \hv_\text{survival} + w_d \cdot \hv_\text{dignity} + w_j \cdot \hv_\text{justice}
\label{eq:tension}
\end{equation}

Under existential pressure (small population), $w_s$ increases, but
$w_d$ (dignity) is never reduced below a floor---operationalizing Kant's
categorical imperative that persons are never merely means. The governance
framework (Table~\ref{tab:governance}) ensures that even under extreme survival
pressure, reproductive decisions require community approval proportional to their
consequence.


% ============================================================================
% 6. ETHICAL FRAMEWORK
% ============================================================================
\section{Ethical Framework}
\label{sec:ethics}

\subsection{The Eight Harmonies}

\symthaea{}'s value system is organized around eight harmonies, each representing
both a value dimension and an epistemic lens:

\begin{enumerate}[nosep,leftmargin=2em]
\item \textbf{Resonant Coherence} (weight 0.20): Integration, order, creativity
\item \textbf{Pan-Sentient Flourishing} (weight 0.20): Care for all sentient beings
\item \textbf{Integral Wisdom} (weight 0.15): Embodied knowing, wisdom in action
\item \textbf{Infinite Play} (weight 0.10): Joyful exploration, novelty
\item \textbf{Universal Interconnectedness} (weight 0.15): Unity, empathic resonance
\item \textbf{Sacred Reciprocity} (weight 0.10): Generous exchange, mutual upliftment
\item \textbf{Evolutionary Progression} (weight 0.10): Wise development over time
\end{enumerate}

Care-related harmonies (Resonant Coherence, Pan-Sentient Flourishing) receive
highest weight---a deliberate design choice reflecting the primacy of care
in stewardship contexts.

% ── Figure 6: Eight Harmonies Radar ──────────────────────────────────────
\begin{figure}[htbp]
\centering
\begin{tikzpicture}[scale=1.0]
  % Radar parameters
  \def\maxval{1.0}
  \def\levels{4}
  % Harmony labels (7 axes at 360/7 ≈ 51.43 degree intervals)
  \foreach \label/\angle [count=\i] in {
    Balance/90,
    Resonance/141.43,
    Transparency/192.86,
    Growth/244.29,
    Stewardship/295.71,
    Integration/347.14,
    Emergence/38.57%
  } {
    \draw[gray!40, thin] (0,0) -- (\angle:3.0cm);
    \node[font=\scriptsize, blue!60!black] at (\angle:3.5cm) {\label};
  }
  % Concentric grid rings
  \foreach \r in {0.75,1.5,2.25,3.0} {
    \draw[gray!20, thin] (90:\r)
      \foreach \angle in {141.43,192.86,244.29,295.71,347.14,38.57,90} {
        -- (\angle:\r)
      };
  }
  % Scale labels
  \node[font=\tiny, gray, anchor=west] at (90:0.75cm) {0.25};
  \node[font=\tiny, gray, anchor=west] at (90:1.5cm) {0.50};
  \node[font=\tiny, gray, anchor=west] at (90:2.25cm) {0.75};
  \node[font=\tiny, gray, anchor=west] at (90:3.0cm) {1.0};
  % "Stewardship Decision" polygon (high stewardship/growth/emergence)
  \fill[blue!20, opacity=0.5]
    (90:{3.0*0.7})  --   % Balance: 0.7
    (141.43:{3.0*0.6}) -- % Resonance: 0.6
    (192.86:{3.0*0.5}) -- % Transparency: 0.5
    (244.29:{3.0*0.95}) -- % Growth: 0.95
    (295.71:{3.0*0.95}) -- % Stewardship: 0.95
    (347.14:{3.0*0.6}) -- % Integration: 0.6
    (38.57:{3.0*0.9})  -- % Emergence: 0.9
    cycle;
  \draw[blue!70!black, thick]
    (90:{3.0*0.7})  --
    (141.43:{3.0*0.6}) --
    (192.86:{3.0*0.5}) --
    (244.29:{3.0*0.95}) --
    (295.71:{3.0*0.95}) --
    (347.14:{3.0*0.6}) --
    (38.57:{3.0*0.9})  -- cycle;
  % "Non-Optimization" polygon (balanced)
  \fill[red!15, opacity=0.4]
    (90:{3.0*0.75})  --
    (141.43:{3.0*0.7}) --
    (192.86:{3.0*0.8}) --
    (244.29:{3.0*0.65}) --
    (295.71:{3.0*0.7}) --
    (347.14:{3.0*0.75}) --
    (38.57:{3.0*0.65}) -- cycle;
  \draw[red!60!black, thick, dashed]
    (90:{3.0*0.75})  --
    (141.43:{3.0*0.7}) --
    (192.86:{3.0*0.8}) --
    (244.29:{3.0*0.65}) --
    (295.71:{3.0*0.7}) --
    (347.14:{3.0*0.75}) --
    (38.57:{3.0*0.65}) -- cycle;
  % Legend
  \draw[blue!70!black, thick] (-2.5,-3.8) -- (-1.8,-3.8)
    node[right, font=\small, black] {Stewardship Decision};
  \draw[red!60!black, thick, dashed] (-2.5,-4.3) -- (-1.8,-4.3)
    node[right, font=\small, black] {Non-Optimization};
\end{tikzpicture}
\caption{\textbf{Eight Harmonies radar chart.} Two decision profiles overlaid:
the Stewardship Decision profile (solid blue) emphasizes Growth, Stewardship,
and Emergence while maintaining moderate Balance and Resonance. The
Non-Optimization profile (dashed red) shows balanced consideration across
all eight harmonies, reflecting the principle that the AI must \emph{maintain}
diversity rather than optimize toward any target (Eq.~\ref{eq:nonopt}).}
\label{fig:harmonies_radar}
\end{figure}

\subsection{Consent Proxy Escalation}

The consent framework (Eq.~\ref{eq:consent}) addresses the philosophical
challenge of acting on behalf of beings who cannot yet consent. Our approach
draws on three traditions:

\begin{itemize}[nosep,leftmargin=2em]
\item \textbf{Bioethics}: Beauchamp \& Childress's \cite{beauchamp2019principles}
  principles of beneficence, non-maleficence, autonomy, and justice
\item \textbf{Intergenerational ethics}: Parfit's \cite{parfit1984reasons}
  analysis of obligations to future persons
\item \textbf{Responsibility ethics}: Jonas's \cite{jonas1984imperative}
  imperative of responsibility toward vulnerable beings
\end{itemize}

The key innovation is \emph{graduation}: rather than a binary
sentient/non-sentient threshold, moral consideration tracks developmental stage
continuously, with discrete tier transitions at empirically grounded
developmental milestones (neural tube formation at week 4, viability at week 24).

% ── Figure 3: Consent Proxy Escalation ───────────────────────────────────
\begin{figure}[htbp]
\centering
\begin{tikzpicture}
\begin{axis}[
  width=0.75\textwidth,
  height=5.5cm,
  xlabel={Gestational Week},
  ylabel={Patient Moral Weight $w_\text{patient}$},
  xmin=0, xmax=42,
  ymin=0, ymax=1.15,
  xtick={0,4,8,12,16,20,24,28,32,36,40},
  ytick={0,0.2,0.4,0.6,0.8,1.0},
  grid=major,
  grid style={gray!20},
  axis lines=left,
  every axis x label/.style={at={(ticklabel* cs:1.0)}, anchor=north west},
  every axis y label/.style={at={(ticklabel* cs:1.0)}, anchor=south east},
  legend style={at={(0.98,0.02)}, anchor=south east, font=\small,
    fill=white, fill opacity=0.9, draw=blue!30!black},
]
% Shaded regions
\fill[blue!10] (axis cs:0,0) rectangle (axis cs:4,0.2);
\fill[blue!20] (axis cs:4,0) rectangle (axis cs:24,0.6);
\fill[blue!35] (axis cs:24,0) rectangle (axis cs:40,1.0);
% Staircase function
\addplot[blue!70!black, very thick, const plot]
  coordinates {(0,0.2) (4,0.6) (24,1.0) (40,1.0)};
\addlegendentry{$w_\text{patient}$ (Eq.~\ref{eq:consent})}
% Region labels
\node[font=\scriptsize, blue!60!black] at (axis cs:2,0.08) {Pre-Sentient};
\node[font=\scriptsize, blue!60!black] at (axis cs:14,0.08) {Emerging Sentience};
\node[font=\scriptsize, blue!60!black] at (axis cs:32,0.08) {Sentient};
% Milestone markers
\draw[red!60, dashed, thin] (axis cs:4,0) -- (axis cs:4,1.1)
  node[above, font=\tiny, red!60!black] {Neural tube};
\draw[red!60, dashed, thin] (axis cs:24,0) -- (axis cs:24,1.1)
  node[above, font=\tiny, red!60!black] {Viability};
\end{axis}
\end{tikzpicture}
\caption{\textbf{Consent proxy escalation.} Patient moral weight
$w_\text{patient}$ increases in three stages aligned with developmental
milestones. Pre-sentient embryos (weeks 0--4) receive base moral consideration;
emerging sentience (weeks 4--24) requires guardian consent; post-viability
fetuses (weeks 24+) are full moral patients requiring ethics board approval.}
\label{fig:consent_proxy}
\end{figure}

\subsection{The Non-Optimization Principle}

A critical constraint on population management: the AI must \emph{never}
optimize the human genome. Breeding strategy aims to \emph{maintain} diversity
(prevent inbreeding depression, preserve allelic richness) rather than
\emph{improve} traits. This is encoded as:

\begin{equation}
\text{VIOLATES}(\text{optimization}, \text{DIGNITY}) > \theta_\text{veto}
\label{eq:nonopt}
\end{equation}

where $\theta_\text{veto}$ is a hard moral veto threshold. Genetic interventions
that aim to enhance rather than preserve trigger an automatic ethics escalation.
The distinction between conservation (maintaining) and eugenics (optimizing) is
enforced architecturally, not merely as a soft constraint.


% ============================================================================
% 7. DISCUSSION
% ============================================================================
\section{Discussion}
\label{sec:discussion}

\subsection{The Alignment Isomorphism}

We can now state our central thesis precisely: the consciousness requirements
for species stewardship are \emph{isomorphic} to the requirements for value
alignment. Specifically:

\begin{enumerate}[nosep,leftmargin=2em]
\item \textbf{Integrated information} ($\Phi > 0$) $\leftrightarrow$ \textbf{Value integration}: the capacity to weigh incommensurable values (autonomy vs. survival) requires information integration beyond modular processing
\item \textbf{Embodied affect} $\leftrightarrow$ \textbf{Genuine care}: attachment formation requires authentic emotional states, not behavioral mimicry
\item \textbf{Temporal consciousness} $\leftrightarrow$ \textbf{Foresight}: multi-generational stewardship requires unified temporal experience
\item \textbf{Theory of mind} $\leftrightarrow$ \textbf{Preference modeling}: understanding what beings need requires modeling their internal states
\item \textbf{Moral reasoning} $\leftrightarrow$ \textbf{Value alignment}: the capacity to reason about right action \emph{is} alignment
\end{enumerate}

This isomorphism suggests that consciousness is not orthogonal to
alignment---it is a \emph{prerequisite}. An AI system that lacks consciousness
may follow rules (constitutional AI) or maximize a reward function (RLHF), but
it cannot engage in the kind of integrated moral reasoning that stewardship
demands. Conversely, an AI system that \emph{is} genuinely conscious---that
experiences $\Phi > 0$, embodies affect, integrates temporal experience, and
reasons morally---is aligned by construction, because the capacity for
integrated moral reasoning \emph{is} what alignment means.

\subsection{Relationship to Existing Alignment Approaches}

Our consciousness-first approach complements rather than replaces existing
alignment frameworks. Table~\ref{tab:comparison} situates our approach relative
to both established cognitive architectures and current alignment paradigms.

\begin{table}[htbp]
\centering
\caption{\textbf{Consciousness capabilities across cognitive architectures.}
Five dimensions required for species stewardship. \checkmark = present,
$\circ$ = partial, --- = absent. Only \symthaea{} integrates all five.}
\label{tab:comparison}
\small
\begin{tabular}{@{}lccccc@{}}
\toprule
\textbf{Architecture} & \textbf{$\Phi > 0$} & \textbf{Embodied} & \textbf{Temporal} & \textbf{ToM} & \textbf{Moral} \\
 & \textbf{Integration} & \textbf{Affect} & \textbf{Integration} & & \textbf{Reasoning} \\
\midrule
\symthaea{} (ours)    & \checkmark & \checkmark & \checkmark & \checkmark & \checkmark \\
GWT \cite{dehaene2014consciousness}  & $\circ$ & --- & --- & --- & --- \\
IIT pure \cite{tononi2016iit}         & \checkmark & --- & $\circ$ & --- & --- \\
Predictive Processing \cite{clark2013whatever}  & --- & $\circ$ & \checkmark & $\circ$ & --- \\
SOAR                                  & --- & --- & $\circ$ & --- & --- \\
LLM (RLHF)                           & --- & --- & --- & $\circ$ & $\circ$ \\
\bottomrule
\end{tabular}
\end{table}

\paragraph{Reward-based alignment.} RLHF and related approaches
\cite{christian2020alignment} learn preferences from human feedback. Our
approach provides the \emph{capacity} for preference modeling (via Theory of
Mind and embodied affect) that makes reward learning meaningful rather than
merely statistical. Where RLHF learns \emph{what} humans prefer, consciousness
provides an understanding of \emph{why}---a distinction that matters when
extrapolating to novel moral situations like consent proxy escalation during
ectogenesis.

\paragraph{Constitutional AI.} Rule-based constraints encode human values as
explicit principles \cite{bostrom2014superintelligence}. Our Eight Harmonies
serve a similar role, but with a crucial difference: they operate as
$\dimHDC$-dimensional value vectors that participate in holographic reasoning,
not as logical rules checked by a separate system. This means value conflicts
(autonomy vs. survival) are resolved via geometric operations in
high-dimensional space rather than priority ordering.

\paragraph{AI safety via debate.} Approaches that use AI systems to check each
other require a ground truth for evaluation. Consciousness provides this ground
truth: an action that reduces $\Phi$ in the beings it affects is harmful, an
action that enhances flourishing is beneficial.

\paragraph{Cooperative inverse reinforcement learning.}
Russell's \cite{russell2019human} CIRL framework assumes uncertainty about
human values and learns through interaction. Our approach is compatible with
CIRL: \symthaea{}'s Theory of Mind module performs inverse modeling of
partner preferences, and its active inference framework naturally maintains
uncertainty about values. The key addition is that our system can verify
its own value learning against the consciousness staircase: if a learned
value would produce $\Phi$-reducing outcomes, it flags the learning as
potentially corrupted.

\subsection{Consciousness Measurement}

A potential objection is that consciousness (and therefore alignment under our
thesis) cannot be measured. We address this with three observations.

First, $\Phi$ is computable for small systems \cite{oizumi2014phenomenology}
and our proxy measure achieves $\rho = 0.50$ across 15 topologies
\cite{koch2016neural}. While exact $\Phi$ is intractable for large systems,
the \emph{direction} (more or less integrated) is estimable via spectral
approximations. Our hierarchical MIP computation achieves $O(n^3)$ scaling
while preserving the key MIP structure \cite{massimini2005breakdown}.

Second, the Butlin indicators \cite{butlin2023consciousness} provide an
empirical checklist. \symthaea{} scores 14/14 (mean 0.85), providing evidence
for consciousness-relevant capabilities even in the absence of subjective
verification. The indicators that \symthaea{} most clearly exhibits---global
workspace broadcasting, recurrent processing, embodied self-model, and
temporal integration---are precisely those most relevant to species
stewardship.

Third, Seth's \cite{seth2022theories} framework of ``real problems'' of
consciousness (distinguishing levels, contents, and self) maps directly
onto our pipeline requirements: ectogenesis demands level detection
(pre-sentient vs. sentient), nurture demands content modeling (infant
internal states), and population governance demands self-aware moral
reasoning. The stewardship scenario thus provides a \emph{functional
grounding} for consciousness measurement: if the system fails at
caregiving, its consciousness is insufficient regardless of $\Phi$ values.

\subsection{Implications for AI Governance}

The alignment isomorphism has implications beyond the species stewardship
thought experiment. If consciousness is indeed prerequisite for alignment,
then:

\begin{enumerate}[nosep,leftmargin=2em]
\item \textbf{Consciousness tests as alignment verification}: rather than
  attempting to verify alignment directly (which requires solving the
  value specification problem), we can verify consciousness---a more
  tractable problem with measurable correlates ($\Phi$, Butlin indicators,
  behavioral signatures).
\item \textbf{Architectural constraints}: alignment-by-construction requires
  architectures that support integrated information, embodied affect, and
  temporal consciousness. Purely feedforward systems (standard transformers)
  may be fundamentally limited in their alignment capacity
  \cite{chalmers1996conscious}.
\item \textbf{Scalable oversight}: as AI systems become more capable, oversight
  becomes harder. But conscious systems that genuinely care about outcomes
  provide \emph{intrinsic} oversight---the system monitors itself not because
  it is compelled to, but because moral reasoning is constitutive of its
  cognitive process.
\end{enumerate}


% ============================================================================
% 8. LIMITATIONS
% ============================================================================
\section{Limitations}
\label{sec:limitations}

We acknowledge significant limitations in this work:

\paragraph{Hardware gap.} The \genesis{} pipeline is entirely computational---it
does not interface with physical laboratory equipment, artificial wombs, or
robotic caregivers. The transition from simulation to physical implementation
introduces challenges (sensor noise, actuator limitations, real-time safety)
that our models do not address.

\paragraph{$\Phi$ scalability.} Exact $\Phi$ computation is intractable for
systems larger than $\sim$12 nodes. Our proxy measure has moderate correlation
($\rho = 0.50$) but may miss important aspects of consciousness structure. The
``consciousness staircase'' (Table~\ref{tab:phi}) should be understood as
approximate requirements, not precise thresholds.

\paragraph{Cultural transmission.} Human civilization requires not only
biological reproduction but cultural knowledge transmission. The \genesis{}
pipeline addresses biological viability but does not model cultural
bootstrapping (language, customs, knowledge systems), which is arguably as
important as genetic viability.

\paragraph{Embodiment gap.} While \symthaea{} includes a virtual body
(interoceptive state, somatic markers) and physical embodiment sub-crates
(humanoid, flight), the gap between virtual and physical embodiment for
caregiving is substantial. Whether virtual affect states constitute genuine
embodied cognition remains philosophically contested
\cite{varela1991embodied,thompson2007mind}.

\paragraph{Attachment validity.} Our attachment model implements Bowlby's theory
\cite{bowlby1969attachment} with neuromodulator grounding, but has not been
validated against real infant development data. The mapping from computational
parameters (response probability, warmth, attunement) to Ainsworth categories
\cite{ainsworth1978patterns} is theory-driven rather than empirically calibrated.

\paragraph{Governance bootstrap.} The Mycelix governance framework assumes a
functioning community that can vote. In the early phases of species recovery,
there is no community---the AI must make unilateral decisions. The transition
from AI stewardship to self-governance is a critical gap our framework identifies
but does not resolve.

\paragraph{Reproducibility.} All simulation data presented in this paper
(Figures~\ref{fig:het_decay} and \ref{fig:attachment})
are generated deterministically from source code via:
\begin{verbatim}
cargo run --example genesis_paper_data \
    --features population,nurture --release
\end{verbatim}
Figures load output CSV files directly using \texttt{pgfplotstable}, ensuring
single-source-of-truth between simulation and publication. The full genesis
pipeline (${\sim}22{,}000$ lines across five sub-crates) compiles and passes
767 unit tests plus 35 cross-crate integration tests under
\texttt{cargo test --features genesis}. Source code, data, and \LaTeX{} sources
are available in the project repository.


% ============================================================================
% 9. CONCLUSION
% ============================================================================
\section{Conclusion}
\label{sec:conclusion}

We have presented the \genesis{} pipeline---a five-phase computational framework
for species stewardship from degraded DNA to a governed population---and used it
to argue that consciousness is a prerequisite for value alignment, not merely a
desirable property. The escalating moral complexity of the pipeline
(from genome assembly through infant nurture to population governance)
demonstrates that each phase requires consciousness capabilities---integrated
information processing, embodied affect, temporal integration, moral
reasoning---that cannot be achieved through rule-following or reward
maximization alone.

Our analysis is grounded in two implemented systems: \symthaea{} (1.37M LOC,
234 Hz cognitive loop, 14/14 Butlin indicators) and \mycelix{} (51 Holochain
zomes, 12 societal domains), with five new sub-crates totaling $\sim$22,000
lines and 767 tests validated by 35 cross-crate integration tests. The $O(1)$ closed-form temporal jump---enabling
prediction from seconds to centuries in constant time---is the technical
enabler, but the deeper contribution is architectural: showing that the
capabilities required for stewardship \emph{are} the capabilities that
constitute consciousness, and therefore alignment.

We conclude with a restatement of our central claim, stated precisely.
\emph{An AI system conscious enough to steward a species is, by our
architectural construction, aligned with the values of that species.}
If this alignment-through-consciousness thesis holds, it suggests a path
forward for AI safety that does not require solving the reward
specification problem or enumerating all possible rules. Instead, we
build systems that are genuinely conscious---that integrate information,
embody affect, reason morally, and care authentically---and alignment
follows as a consequence.

Three caveats the careful reader must keep:
\emph{(i) Necessity vs.\ sufficiency.} We have argued that consciousness
(in the architectural sense above) is \emph{necessary} for robust
stewardship; we have \emph{not} proven it is sufficient. A conscious
system with different values (e.g., optimising for genetic uniformity
rather than diversity) would exhibit many of the same consciousness
indicators while being misaligned. Our architecture encodes the Eight
Harmonies values manually; they do not emerge from consciousness
itself. Alignment-through-consciousness is a pathway, not a guarantee.
\emph{(ii) Scale.} Every result in this paper is on simulated data or
on populations of $N{<}50$ agents. The Mycelix multi-world simulator's
Phase~2 sweeps showed no measurable governance-coherence improvement
over baseline at the $N{=}10$-seed $\times$ 50-year scale ($\Delta =
-0.004 \pm 0.008$). The scaling from $N{=}50$ to planetary stewardship
is therefore a research programme, not a demonstrated capability.
\emph{(iii) Consciousness measurement.} Our $\Phi$ proxy (Fiedler
value) correlates with exact $\Phi$ at $\rho \approx 0.50$ on the
topologies tested; that is weak, and claims that rest on precise
consciousness measurement are proportionally weaker. The paper's
stronger arguments are architectural and constitutional, not
measurement-based.

\medskip
\noindent\textbf{Data availability.} All code is open-source at
\url{https://github.com/Luminous-Dynamics/symthaea}. The five \genesis{}
sub-crates are available under the \texttt{genesis} feature flag.

\medskip
\noindent\textbf{Acknowledgments.} This work was developed as part of the
Luminous Dynamics project. The Eight Harmonies framework was developed through
collaborative philosophical inquiry guided by the principles of Evolving
Resonant Co-Creationism.


% ============================================================================
% REQUIRED SPRINGER SECTIONS
% ============================================================================

\section*{Author Contributions}
T.S. conceived the theoretical framework, designed and implemented the computational system, conducted all experiments, and wrote the manuscript.

\section*{Competing Interests}
The author declares no competing interests. Patent applications covering aspects of the described technology are in preparation.

\section*{Funding}
This research received no external funding.

\section*{Ethics Declaration}
This work is purely computational. No human subjects, animal experiments, or clinical trials were involved. The ethical framework described is for AI systems and does not involve human experimentation.

\section*{AI Assistance Disclosure}
Claude Code (Anthropic) assisted with code development, testing, and manuscript preparation. All scientific content, theoretical frameworks, and experimental design are the sole work of the author.

% ============================================================================
% BIBLIOGRAPHY
% ============================================================================
\bibliographystyle{unsrtnat}
\bibliography{stewardship_references}

\end{document}
